\documentclass[10pt]{article}
\usepackage{amsmath,amssymb}
\usepackage{longtable}
\usepackage{xcolor}
\usepackage[a4paper,left=2cm,right=2cm,top=1.5cm,bottom=2cm]{geometry}
\setlength{\emergencystretch}{2em}

\setlength{\parindent}{0pt}

\begin{document}

\title{Forward and Reverse Automatic Differentiation by Hand - Solutions}
\author{}
\date{}
\maketitle
\subsection*{Learning goals}
After this exercise, you should be able to

\begin{itemize}
\item
explain what a \textbf{\textit{dual number}} is and why dual numbers are useful for forward-mode automatic differentiation,

\item
propagate values and derivatives through a computation using dual numbers,

\item
compute directional derivatives and Jacobian columns with forward mode,

\item
explain the idea of a \textbf{\textit{computational graph}} and an \textbf{\textit{adjoint}} in reverse-mode automatic differentiation,

\item
propagate adjoints backwards through a computation graph,

\item
compute gradients and Jacobian rows with reverse mode,

\item
explain the practical difference between forward and reverse AD.

\end{itemize}
The calculations in this sheet are designed to be done \textbf{\textit{by hand}}. No programming is required.

\clearpage
\section*{Part I - Forward-mode AD with dual numbers}
\subsection*{1. What are dual numbers?}
Automatic differentiation (AD) computes derivatives by applying the chain rule to the elementary operations of a calculation. For a composition, the chain rule gives $(f\circ q)'(x)=f'(q(x))q'(x)$.

A \textbf{\textit{dual number}} $\hat{x}$ has the form

\[
\hat{x} = x + \dot{x}\varepsilon
\]

where $x,\dot{x}\in\mathbb{R}$ and $\varepsilon$ is a formal symbol with the special property

\[
\varepsilon^2 = 0
\qquad
\varepsilon \neq 0
\]

The ordinary value $x$ is called the \textbf{\textit{primal value}}. The coefficient $\dot{x}$ is the \textbf{\textit{tangent}}: it tracks the derivative along a chosen input direction. The dot is notation for this derivative component, not necessarily a time derivative. Choosing the initial tangent is called setting an \textbf{\textit{input seed}}. \textbf{\textit{Forward mode}} carries primal values and tangents together from inputs to outputs.

If we then evaluate a differentiable function $f$ using its value and first derivative for each elementary operation together with $\varepsilon^2=0$, we obtain

\[
f(x+\dot{x}\varepsilon)
= f(x) + f'(x)\dot{x}\,\varepsilon
\]

For the special choice $\dot{x}=1$

\[
f(x+\varepsilon)=f(x)+f'(x)\varepsilon
\]

so the coefficient of $\varepsilon$ is directly $f'(x)$.

\subsection*{2. Arithmetic with dual numbers}
Let

\[
\hat{a}=a+\dot{a}\varepsilon,
\qquad
\hat{b}=b+\dot{b}\varepsilon.
\]

Then

\subsubsection*{Addition}
\[
\hat{a}+\hat{b}
=(a+b)+(\dot{a}+\dot{b})\varepsilon.
\]

\subsubsection*{Multiplication}
\[
\hat{a}\hat{b}
=(a+\dot{a}\varepsilon)(b+\dot{b}\varepsilon).
\]

Expanding the product gives

\[
ab+(a\dot{b}+\dot{a}b)\varepsilon
\]

which is exactly the product rule.

\subsubsection*{Some useful elementary functions}
Here $\log$ denotes the natural logarithm.

\[
\exp(a+\dot{a}\varepsilon)
= e^a+\dot{a}\,e^a\varepsilon,
\]

\[
\sin(a+\dot{a}\varepsilon)
= \sin(a)+\dot{a}\,\cos(a)\varepsilon,
\]

\[
\log(a+\dot{a}\varepsilon)
= \log(a)+\frac{\dot{a}}{a}\varepsilon,
\qquad a>0,
\]

and

\[
(a+\dot{a}\varepsilon)^2
= a^2+2a\,\dot{a}\varepsilon.
\]

\subsection*{3. Very simple 1D example}
Consider

\[
f(x)=x^2+3x
\]

and compute $f(2)$ and $f'(2)$ using dual numbers.

The dual $\hat{x}_0$ is (with tangent $\dot{x}_0=1$)

\[
\hat{x}_0=2+\varepsilon.
\]

Then

\[
\hat{x}_0^2=(2+\varepsilon)^2
=4+4\varepsilon,
\]

and

\[
3\hat{x}_0=6+3\varepsilon.
\]

Therefore

\[
f(\hat{x}_0)
=10+7\varepsilon.
\]

Hence

\[
f(2)=10
\]
\[
f'(2)=7.
\]

The important point is that the value and the derivative were propagated at the same time.

\subsection*{4. More than one input}
A scalar is a single number; a vector is an ordered collection of numbers. Vectors are written as column vectors, and $T$ denotes transposition.

Suppose now that

\[
g:\mathbb{R}^n\to\mathbb{R}^m.
\]

We can seed every input with its own tangent component:

\[
\hat{\mathbf{x}}
=\mathbf{x}+\dot{\mathbf{x}}\varepsilon.
\]

Forward AD then gives

\[
{
 g(\mathbf{x}+\dot{\mathbf{x}}\varepsilon)
 = g(\mathbf{x}) + J_g(\mathbf{x})\,\dot{\mathbf{x}}\,\varepsilon
}
\]

The \textbf{\textit{Jacobian matrix}} $J_g(\mathbf{x})$ collects all first partial derivatives: its entry in row $i$ and column $j$ is $\partial g_i/\partial x_j$. A \textbf{\textit{partial derivative}} measures the change with respect to one input while holding the others fixed. For $n$ inputs and $m$ outputs, the Jacobian has $m$ rows and $n$ columns.

A \textbf{\textit{sweep}} is one traversal of the calculation in the chosen direction. One forward-mode sweep computes one \textbf{\textit{Jacobian-vector product (JVP)}}

\[
J_g(\mathbf{x})\dot{\mathbf{x}}.
\]

A \textbf{\textit{standard basis vector}} has one entry equal to $1$ and all other entries equal to $0$. If we choose $\dot{\mathbf{x}}$ to be a standard basis vector, for example

\[
\mathbf{e}_1=(1,0,0)^T,
\]

then the tangent part is the first column of the Jacobian.

For a scalar-valued function, the \textbf{\textit{gradient}} $\nabla g$ is the column vector of all partial derivatives. The \textbf{\textit{directional derivative}} along an input direction $\mathbf v$ is

\[
D_{\mathbf v}g(\mathbf x)=\left.\frac{d}{dt}g(\mathbf x+t\mathbf v)\right|_{t=0}=\nabla g(\mathbf x)^T\mathbf v.
\]

Here $\mathbf v$ need not have unit length: scaling the seed scales the directional derivative.

\section*{Forward-mode exercises}
\subsection*{Exercise 1.1 - Scalar input, scalar output}
Consider

\[
f(x)=\log\left(xe^x+3\right).
\]

Evaluate the function and its derivative at

\[
x=1
\]

using dual-number propagation.

\begingroup\color{red}
\paragraph*{Solution 1.1}
We start with

\[
\hat{x}=1+\varepsilon.
\]

\subsubsection*{Step 1: exponential}
\[
a=e^{\hat{x}}
=e^{1+\varepsilon}
=e+e\varepsilon.
\]

\subsubsection*{Step 2: multiplication}
\[
b=\hat{x}a
=(1+\varepsilon)(e+e\varepsilon).
\]

Expanding and dropping $\varepsilon^2$,

\[
b=e+2e\varepsilon.
\]

\subsubsection*{Step 3: addition}
\[
c=b+3
=(e+3)+2e\varepsilon.
\]

\subsubsection*{Step 4: logarithm}
Using

\[
\log(a+b\varepsilon)
=\log a+\frac{b}{a}\varepsilon,
\]

we obtain

\[
f(\hat{x})
=
\log(e+3)
+
\frac{2e}{e+3}\varepsilon.
\]

Therefore

\[
{f(1)=\log(e+3)}
\]

and

\[
{
f'(1)=\frac{2e}{e+3}
}.
\]
\par\endgroup

\subsection*{Exercise 1.2 - Vector input, scalar output}
Consider

\[
g(x,y,z)=xy+\sin z+y^2
\]

at the point

\[
(x,y,z)=(1,2,0).
\]

\begin{enumerate}
\item
First calculate the ordinary function value $g(1,2,0)$.

\item
Use forward-mode AD with the three seeds

\[
\mathbf{e}_x=(1,0,0),
   \qquad
   \mathbf{e}_y=(0,1,0),
   \qquad
   \mathbf{e}_z=(0,0,1)
\]

to determine the full gradient

\[
\nabla g(1,2,0).
\]

\item
Now use only one forward sweep with seed direction

\[
\dot{\mathbf{x}}=(1,-1,2)^T
\]

to compute the directional derivative

\[
D_{\dot{\mathbf{x}}}g
   =
   \nabla g^T\dot{\mathbf{x}}.
\]

\end{enumerate}

\begingroup\color{red}
\paragraph*{Solution 1.2}
The function is

\[
g(x,y,z)=xy+\sin z+y^2.
\]

At $(1,2,0)$,

\[
g(1,2,0)=1\cdot2+\sin 0+2^2=6.
\]

\subsubsection*{Seed in the $x$-direction}
Use

\[
\hat{x}=1+\varepsilon,
\qquad
\hat{y}=2,
\qquad
\hat{z}=0.
\]

Then

\[
\hat{x}\hat{y}
=(1+\varepsilon)2
=2+2\varepsilon,
\]

\[
\sin\hat{z}=0,
\]

and

\[
\hat{y}^2=4.
\]

Thus

\[
g=6+2\varepsilon.
\]

Therefore

\[
\frac{\partial g}{\partial x}=2.
\]

\subsubsection*{Seed in the $y$-direction}
Use

\[
\hat{x}=1,
\qquad
\hat{y}=2+\varepsilon,
\qquad
\hat{z}=0.
\]

Then

\[
\hat{x}\hat{y}=2+\varepsilon,
\]

and

\[
\hat{y}^2
=(2+\varepsilon)^2
=4+4\varepsilon.
\]

Hence

\[
g=6+5\varepsilon,
\]

so

\[
\frac{\partial g}{\partial y}=5.
\]

\subsubsection*{Seed in the $z$-direction}
Use

\[
\hat{z}=\varepsilon.
\]

Since

\[
\sin(\varepsilon)
=\sin 0+\cos 0\,\varepsilon
=\varepsilon,
\]

we get

\[
g=6+\varepsilon.
\]

Therefore

\[
\frac{\partial g}{\partial z}=1.
\]

The gradient is

\[
{
\nabla g(1,2,0)
=
\begin{pmatrix}
2\\5\\1
\end{pmatrix}
}.
\]

\subsubsection*{Directional derivative}
For

\[
\mathbf v=(1,-1,2)^T,
\]

seed all inputs at once:

\[
\hat{x}=1+\varepsilon,
\qquad
\hat{y}=2-\varepsilon,
\qquad
\hat{z}=2\varepsilon.
\]

Forward propagation gives

\[
\hat x\hat y=2+\varepsilon,\qquad
\sin\hat z=2\varepsilon,\qquad
\hat y^2=4-4\varepsilon.
\]

Adding the three terms gives $g=6-\varepsilon$. The tangent is therefore $-1$, in agreement with $2(1)+5(-1)+1(2)=-1$.

Thus

\[
{D_{\mathbf v}g=-1}.
\]
\par\endgroup

\subsection*{Exercise 1.3 - Vector input, vector output}
Consider the vector-valued function

\[
\mathbf{h}(x,y,z)
=
\begin{pmatrix}
 h_1(x,y,z)\\
 h_2(x,y,z)
\end{pmatrix}
=
\begin{pmatrix}
 xy+z\\
 x^2+\sin y-z^2
\end{pmatrix}.
\]

Evaluate it at

\[
(x,y,z)=(1,0,1).
\]

\begin{enumerate}
\item
Calculate $\mathbf h(1,0,1)$.

\item
Use the three input basis directions

\[
\mathbf e_x,
   \qquad
   \mathbf e_y,
   \qquad
   \mathbf e_z
\]

and dual numbers to compute the full Jacobian

\[
J_{\mathbf h}
   =
   \begin{pmatrix}
   \frac{\partial h_1}{\partial x} &
   \frac{\partial h_1}{\partial y} &
   \frac{\partial h_1}{\partial z}\\[4pt]
   \frac{\partial h_2}{\partial x} &
   \frac{\partial h_2}{\partial y} &
   \frac{\partial h_2}{\partial z}
   \end{pmatrix}.
\]

Remember: each forward sweep gives one column of the Jacobian.

\item
Use a single forward sweep with

\[
\mathbf v=(1,2,-1)^T
\]

to compute

\[
J_{\mathbf h}\mathbf v
\]

\end{enumerate}

\begingroup\color{red}
\paragraph*{Solution 1.3}
At $(1,0,1)$,

\[
h_1=1\cdot0+1=1,
\]

\[
h_2=1^2+\sin 0-1^2=0.
\]

Therefore

\[
{
\mathbf h(1,0,1)
=
\begin{pmatrix}1\\0\end{pmatrix}
}.
\]

\subsubsection*{Seed $\mathbf e_x$}
Use

\[
\hat{x}=1+\varepsilon,
\qquad
\hat{y}=0,
\qquad
\hat{z}=1.
\]

For the first component,

\[
h_1=(1+\varepsilon)0+1=1+0\varepsilon.
\]

For the second component,

\[
h_2=(1+\varepsilon)^2+\sin 0-1
=2\varepsilon.
\]

So the first Jacobian column is

\[
\begin{pmatrix}0\\2\end{pmatrix}.
\]

\subsubsection*{Seed $\mathbf e_y$}
Use

\[
\hat{y}=\varepsilon.
\]

Then

\[
h_1=1+\varepsilon,
\]

and

\[
h_2=\sin(\varepsilon)=\varepsilon.
\]

So the second Jacobian column is

\[
\begin{pmatrix}1\\1\end{pmatrix}.
\]

\subsubsection*{Seed $\mathbf e_z$}
Use

\[
\hat{z}=1+\varepsilon.
\]

Then

\[
h_1=1+\varepsilon,
\]

while

\[
h_2
=1-(1+\varepsilon)^2
=-2\varepsilon.
\]

So the third Jacobian column is

\[
\begin{pmatrix}1\\-2\end{pmatrix}.
\]

Therefore

\[
{
J_{\mathbf h}(1,0,1)
=
\begin{pmatrix}
0 & 1 & 1\\
2 & 1 & -2
\end{pmatrix}
}.
\]

For

\[
\mathbf v=(1,2,-1)^T,
\]

seed the inputs as $\hat x=1+\varepsilon$, $\hat y=2\varepsilon$, and $\hat z=1-\varepsilon$. One forward sweep gives

\[
\hat h_1=(1+\varepsilon)2\varepsilon+(1-\varepsilon)=1+\varepsilon,
\]

\[
\hat h_2=(1+\varepsilon)^2+\sin(2\varepsilon)-(1-\varepsilon)^2=6\varepsilon.
\]

The tangent vector is $(1,6)^T$, agreeing with multiplication of the Jacobian by $\mathbf v$.

Thus

\[
{
J_{\mathbf h}\mathbf v
=
\begin{pmatrix}1\\6\end{pmatrix}
}.
\]
\par\endgroup

\clearpage
\section*{Part II - Reverse-mode AD / Backward AD}
\subsection*{1. The basic idea}
Forward mode propagates derivatives together with the values from input to output.

A \textbf{\textit{computational graph}} represents a calculation as nodes for inputs and intermediate results, with directed edges showing which values each operation uses. An \textbf{\textit{intermediate variable}} stores the result of one of these operations. A \textbf{\textit{local derivative}} is the derivative of an operation's output with respect to one of its inputs.

\textbf{\textit{Reverse mode}} proceeds as follows:

\begin{enumerate}
\item
Perform an ordinary forward pass and store the intermediate values.

\item
Start from the output.

\item
Propagate sensitivities backwards through the computational graph.

\end{enumerate}
For an intermediate variable $v$, define its \textbf{\textit{adjoint}} (the sensitivity of the chosen scalar output to this variable) as

\[
{
\bar v = \frac{\partial L}{\partial v}
}
\]

where $L$ is the final scalar output whose derivative we want (the bar notation is common in reverse-mode AD).

Initialize all adjoints to zero. The \textbf{\textit{output seed}} specifies the starting adjoint at the output. For the derivative of the scalar output itself, set

\[
\bar L
=
\frac{\partial L}{\partial L}
=1.
\]

Then the chain rule is applied locally, one operation at a time, in reverse order.

\subsection*{2. Local backward rules}
Suppose an intermediate variable $c$ is computed from earlier variables.

\subsubsection*{Addition}
If

\[
c=a+b,
\]

then

\[
\bar a \mathrel{+}= \bar c,
\qquad
\bar b \mathrel{+}= \bar c.
\]

\textbf{\textit{Accumulation}}, written $\bar a\mathrel{+}=\bar c$, means replacing $\bar a$ by its current value plus $\bar c$: if one variable influences the output through several paths, all contributions to its adjoint must be added.

\subsubsection*{Multiplication}
If

\[
c=ab,
\]

then

\[
\bar a \mathrel{+}= \bar c\,b,
\qquad
\bar b \mathrel{+}= \bar c\,a.
\]

\subsubsection*{Some useful elementary functions}
\[
c = e^a \qquad \Rightarrow \qquad \bar a \mathrel{+}= \bar c\,e^a
\]
\[
c=\sin a \qquad \Rightarrow \qquad \bar a \mathrel{+}= \bar c\cos a.
\]
\[
c=\log a \qquad \Rightarrow \qquad \bar a \mathrel{+}= \bar c\frac{1}{a}.
\]
\[
c=a^2 \qquad \Rightarrow \qquad \bar a \mathrel{+}= \bar c\,2a.
\]

\subsection*{3. Very simple 1D reverse-mode example}
Consider

\[
f(x)=(x+1)^2
\]

at $x=2$.

Break the function into elementary operations:

\[
v_1=x+1,
\]

\[
v_2=v_1^2,
\]

\[
f=v_2.
\]

\subsubsection*{Forward pass}
At $x=2$,

\[
v_1=3,
\qquad
v_2=9.
\]

\subsubsection*{Backward pass}
Start with

\[
\bar v_2=1
\]
because $f=v_2$, so $\partial f/\partial v_2=1$.

Since

\[
v_2=v_1^2
\]

we get

\[
\bar v_1
=
\bar v_2\,2v_1
=1\cdot 6
=6
\]

and since

\[
v_1=x+1
\]

we obtain

\[
\bar x=\bar v_1=6
\]

Hence

\[
f'(2)=6.
\]

In one dimension this may look more complicated than ordinary differentiation. Its advantage becomes clear when a function has many inputs but only one scalar output.

\subsection*{4. Reverse mode for several inputs}
For

\[
f:\mathbb R^n\to\mathbb R,
\]

one reverse sweep gives

\[
\nabla f
\]

with derivatives with respect to all inputs at once.

A loss measures the error of a model's predictions. In machine learning it is usually a scalar, even when the model has many parameters.

For a vector-valued function

\[
\mathbf f:\mathbb R^n\to\mathbb R^m,
\]

we must first choose an output seed $\mathbf w\in\mathbb R^m$. Reverse mode then computes

\[
{
J_{\mathbf f}^T\mathbf w
}.
\]

The \textbf{\textit{vector-Jacobian product (VJP)}} is $\mathbf w^TJ_{\mathbf f}$. With column-vector notation, we write its transpose $J_{\mathbf f}^T\mathbf w$. The output seed $\mathbf w$ assigns a starting adjoint to each output; equivalently, we differentiate the scalar $L=\mathbf w^T\mathbf f$.

To reconstruct the complete Jacobian, use one reverse sweep per output basis vector.

\section*{Reverse-mode exercises}
Use exactly the same three functions as in the forward-mode section.

\subsection*{Exercise 2.1 - Scalar input, scalar output}
For

\[
f(x)=\log(xe^x+3)
\]

at $x=1$:

\begin{enumerate}
\item
Write the function as a sequence of elementary intermediate variables.

\item
Perform the forward pass and record all intermediate values.

\item
Initialize all adjoints to zero, then set the output adjoint equal to $1$.

\item
Propagate all adjoints backwards.

\item
Determine $\bar x=f'(1)$.

\end{enumerate}

\begingroup\color{red}
\paragraph*{Solution 2.1}
Define

\[
v_1=e^x,
\qquad
v_2=xv_1,
\qquad
v_3=v_2+3,
\qquad
v_4=\log v_3.
\]

The final output is

\[
f=v_4.
\]

\subsubsection*{Forward pass}
At $x=1$,

\[
v_1=e,
\]

\[
v_2=e,
\]

\[
v_3=e+3,
\]

\[
v_4=\log(e+3).
\]

\subsubsection*{Backward pass}
Initialize all adjoints to zero before setting the output seed.

Initialize

\[
\bar v_4=1.
\]

Because

\[
v_4=\log v_3,
\]

\[
\bar v_3
=
\bar v_4\frac{1}{v_3}
=
\frac{1}{e+3}.
\]

Because

\[
v_3=v_2+3,
\]

\[
\bar v_2
=
\bar v_3
=
\frac{1}{e+3}.
\]

Now

\[
v_2=xv_1.
\]

This gives one direct contribution to $x$:

\[
\bar x_{\text{direct}}
=
\bar v_2 v_1
=
\frac{e}{e+3}.
\]

It also gives

\[
\bar v_1
=
\bar v_2 x
=
\frac{1}{e+3}.
\]

Finally,

\[
v_1=e^x,
\]

so the indirect contribution to $x$ is

\[
\bar x_{\text{via }v_1}
=
\bar v_1 e^x
=
\frac{e}{e+3}.
\]

The two paths must be added:

\[
\bar x
=
\frac{e}{e+3}
+
\frac{e}{e+3}.
\]

Therefore

\[
{
f'(1)=\bar x=\frac{2e}{e+3}
}.
\]

This is exactly the same derivative obtained with dual numbers.
\par\endgroup

\subsection*{Exercise 2.2 - Three inputs, scalar output}
For

\[
g(x,y,z)=xy+\sin z+y^2
\]

at

\[
(x,y,z)=(1,2,0),
\]

use the intermediate variables

\[
a=xy,
\qquad
b=\sin z,
\qquad
c=y^2,
\]

\[
d=a+b+c
\]

The scalar output is $g=d$.

\begin{enumerate}
\item
Perform the forward pass and record $a,b,c,d$.

\item
Initialize all adjoints to zero, then set $\bar d=1$ and propagate backwards.

\item
Determine the gradient $\nabla g(1,2,0)$ and compare it with Exercise 1.2.

\item
Explain why the adjoint of $y$ receives two contributions.

\end{enumerate}

\begingroup\color{red}
\paragraph*{Solution 2.2}
Use

\[
a=xy,
\qquad
b=\sin z,
\qquad
c=y^2,
\qquad
d=a+b+c
\]

The scalar output is $g=d$.

\subsubsection*{Forward pass}
At $(1,2,0)$,

\[
a=2,
\qquad
b=0,
\qquad
c=4,
\qquad
d=6
\]

\subsubsection*{Backward pass}
Initialize all adjoints to zero before setting the output seed.

Start with

\[
\bar d=1.
\]

Since

\[
d=a+b+c,
\]

we obtain

\[
\bar a=1,
\qquad
\bar b=1,
\qquad
\bar c=1
\]

From

\[
a=xy,
\]

we get

\[
\bar x
\mathrel{+}=
\bar a\,y
=1\cdot2
=2,
\]

and

\[
\bar y
\mathrel{+}=
\bar a\,x
=1\cdot1
=1.
\]

From

\[
b=\sin z,
\]

we obtain

\[
\bar z
\mathrel{+}=
\bar b\cos z
=1\cdot1
=1.
\]

From

\[
c=y^2,
\]

we obtain another contribution to $y$:

\[
\bar y
\mathrel{+}=
\bar c\,2y
=1\cdot4
=4.
\]

Therefore

\[
\bar y=1+4=5.
\]

The final gradient is

\[
{
\nabla g(1,2,0)
=
\begin{pmatrix}
2\\5\\1
\end{pmatrix}
}.
\]

The adjoint of $y$ gets two contributions because $y$ influences the output through both

\[
xy
\]

and

\[
y^2.
\]

Reverse mode must accumulate both paths.

This agrees with the gradient obtained in Exercise 1.2.
\par\endgroup

\subsection*{Exercise 2.3 - Three inputs, two outputs}
For

\[
\mathbf h(x,y,z) = \begin{pmatrix} xy+z\\ x^2+\sin y-z^2\end{pmatrix}
\]

at

\[
(x,y,z)=(1,0,1),
\]

reverse mode needs an output seed because the output is not scalar.

\begin{enumerate}
\item
First output. Use

\[
\mathbf w_1 = \begin{pmatrix}1\\0\end{pmatrix}.
\]

That is, propagate backwards only from $h_1$.

Compute

\[
J_{\mathbf h}^T\mathbf w_1.
\]

\item
Second output. Use

\[
\mathbf w_2 = \begin{pmatrix}0\\1\end{pmatrix}
\]

Compute

\[
J_{\mathbf h}^T\mathbf w_2
\]

\item
Full Jacobian. Use the two results to reconstruct the complete Jacobian.

\item
General output seed. Without calculating the Jacobian from scratch, use

\[
\mathbf w = \begin{pmatrix}3\\-1\end{pmatrix}
\]

to calculate

\[
J_{\mathbf h}^T\mathbf w
\]

\end{enumerate}

\begingroup\color{red}
\paragraph*{Solution 2.3}
The two outputs are

\[
h_1=xy+z,
\]

and

\[
h_2=x^2+\sin y-z^2.
\]

At $(1,0,1)$,

\[
\mathbf h(1,0,1)
=
\begin{pmatrix}1\\0\end{pmatrix}.
\]

\subsubsection*{Reverse sweep from $h_1$}
Use output seed

\[
\mathbf w_1
=
\begin{pmatrix}1\\0\end{pmatrix}.
\]

Only $h_1$ contributes.

Since

\[
h_1=xy+z,
\]

we have

\[
\frac{\partial h_1}{\partial x}=y=0,
\]

\[
\frac{\partial h_1}{\partial y}=x=1,
\]

\[
\frac{\partial h_1}{\partial z}=1.
\]

Thus

\[
{
J_{\mathbf h}^T\mathbf w_1
=
\begin{pmatrix}
0\\1\\1
\end{pmatrix}
}.
\]

This is the first \textbf{\textit{row}} of the Jacobian, written as a column vector.

\subsubsection*{Reverse sweep from $h_2$}
Use

\[
\mathbf w_2
=
\begin{pmatrix}0\\1\end{pmatrix}.
\]

Since

\[
h_2=x^2+\sin y-z^2,
\]

we obtain

\[
\frac{\partial h_2}{\partial x}=2x=2,
\]

\[
\frac{\partial h_2}{\partial y}=\cos y=1,
\]

\[
\frac{\partial h_2}{\partial z}=-2z=-2.
\]

Therefore

\[
{
J_{\mathbf h}^T\mathbf w_2
=
\begin{pmatrix}
2\\1\\-2
\end{pmatrix}
}.
\]

Putting the two rows together gives

\[
{
J_{\mathbf h}(1,0,1)
=
\begin{pmatrix}
0 & 1 & 1\\
2 & 1 & -2
\end{pmatrix}
}.
\]

\subsubsection*{General output seed}
Let

\[
\mathbf w
=
\begin{pmatrix}3\\-1\end{pmatrix}.
\]

Then

\[
J_{\mathbf h}^T\mathbf w
=
3\nabla h_1-\nabla h_2.
\]

Therefore

\[
J_{\mathbf h}^T\mathbf w
=
3
\begin{pmatrix}0\\1\\1\end{pmatrix}
-
\begin{pmatrix}2\\1\\-2\end{pmatrix}
=
\begin{pmatrix}-2\\2\\5\end{pmatrix}.
\]

Hence

\[
{
J_{\mathbf h}^T\mathbf w
=
\begin{pmatrix}-2\\2\\5\end{pmatrix}
}.
\]

\subsubsection*{Local reverse-mode propagation}
Use intermediate variables
\[
a=xy,\quad b=x^2,\quad c=\sin y,\quad d=z^2,\quad q=b+c,
\qquad h_1=a+z,\quad h_2=q-d.
\]
At $(1,0,1)$ their values are $(a,b,c,d,q)=(0,1,0,1,1)$. For output seed $(w_1,w_2)^T$, initialize all other adjoints to zero and set $\bar h_1=w_1$, $\bar h_2=w_2$.
The output operations give
\[
\bar a=w_1,\quad \bar z=w_1,\quad \bar q=w_2,\quad \bar d=-w_2,
\qquad \bar b=\bar c=w_2.
\]
Propagating through the remaining operations and accumulating at the inputs gives
\[
\bar x=\bar a\,y+2x\bar b=2w_2,\qquad
\bar y=\bar a\,x+\bar c\cos y=w_1+w_2,\qquad
\bar z=w_1+2z\bar d=w_1-2w_2.
\]
Thus the seeds $(1,0)^T$, $(0,1)^T$ and $(3,-1)^T$ give $(0,1,1)^T$, $(2,1,-2)^T$ and $(-2,2,5)^T$, respectively.
\par\endgroup

\clearpage
\section*{Part III - Short conceptual questions}
\subsection*{Question 1}
Why does setting $\varepsilon^2=0$ cause the coefficient of $\varepsilon$ to behave like a derivative?

\begingroup\color{red}
\paragraph*{Solution 1}
When arithmetic is expanded with

\[
\varepsilon^2=0,
\]

all terms of second and higher order in $\varepsilon$ vanish. What remains is exactly the first-order term of the function expansion:

\[
f(x+\dot x\varepsilon)
=f(x)+f'(x)\dot x\varepsilon.
\]

Therefore the coefficient of $\varepsilon$ follows the chain rule automatically.
\par\endgroup

\subsection*{Question 2}
For a function

\[
f:\mathbb R^{100}\to\mathbb R,
\]

how many basis-direction forward sweeps would be needed to obtain the full gradient? How many reverse sweeps?

\begingroup\color{red}
\paragraph*{Solution 2}
For

\[
f:\mathbb R^{100}\to\mathbb R,
\]

the full gradient has 100 input derivatives.

Forward mode needs 100 basis-direction sweeps.

Reverse mode needs only one reverse sweep because the output is scalar.
\par\endgroup

\subsection*{Question 3}
For a function

\[
f:\mathbb R\to\mathbb R^{100},
\]

which mode would naturally be more efficient for constructing the full Jacobian?

\begingroup\color{red}
\paragraph*{Solution 3}
For

\[
f:\mathbb R\to\mathbb R^{100},
\]

there is only one input direction. One forward sweep gives the complete single Jacobian column, i.e. all 100 output derivatives with respect to the one input.

Forward mode is therefore the natural choice.
\par\endgroup

\subsection*{Question 4}
What is an adjoint $\bar v$?

\begingroup\color{red}
\paragraph*{Solution 4}
The adjoint of an intermediate variable $v$ is

\[
\bar v=\frac{\partial L}{\partial v},
\]

where $L$ is the final scalar output being differentiated.

It measures how sensitive the final output is to a change in that intermediate variable.
\par\endgroup

\subsection*{Question 5}
Why do reverse-mode updates use accumulation such as

\[
\bar x\mathrel{+}=\cdots
\]

instead of simply assigning one value to $\bar x$?

\begingroup\color{red}
\paragraph*{Solution 5}
A variable may influence the output through several different paths in the computational graph. Each path contributes to the total derivative. Reverse mode must therefore add all contributions to the variable's adjoint.
\par\endgroup

\begingroup\color{red}
\subsection*{Forward mode versus reverse mode}
The most important distinction is not that one method is "better" than the other. They compute different Jacobian products efficiently.

For

\[
f:\mathbb R^n\to\mathbb R^m
\]

with Jacobian

\[
J\in\mathbb R^{m\times n},
\]

we have:

\begingroup
\small
\setlength{\tabcolsep}{4pt}
\begin{longtable}{p{\dimexpr 0.225\textwidth-2\tabcolsep\relax}p{\dimexpr 0.225\textwidth-2\tabcolsep\relax}p{\dimexpr 0.225\textwidth-2\tabcolsep\relax}p{\dimexpr 0.225\textwidth-2\tabcolsep\relax}}
\hline
Mode & Seed & One sweep computes & Full Jacobian requires \\
\hline
\endfirsthead
Mode & Seed & One sweep computes & Full Jacobian requires \\
\hline
\endhead
Forward mode & input direction $\mathbf v\in\mathbb R^n$ & $J\mathbf v$ & $n$ basis sweeps \\[1.5pt]
Reverse mode & output direction $\mathbf w\in\mathbb R^m$ & $J^T\mathbf w$ & $m$ basis sweeps \\[1.5pt]
\hline
\end{longtable}
\endgroup
This gives an important rule of thumb:

\begin{itemize}
\item
\textbf{\textit{Few inputs, many outputs:}} forward mode is often attractive.

\item
\textbf{\textit{Many inputs, few outputs:}} reverse mode is often attractive.

\item
\textbf{\textit{Many parameters, one scalar loss:}} reverse mode is especially attractive.

\end{itemize}
That last case is exactly the situation encountered when training neural networks.

A reverse sweep requires a preceding forward pass to obtain the intermediate values. The sweep counts above refer to one seed direction per sweep.
\par\endgroup

\subsubsection*{Key concepts}
\begingroup
\small
\setlength{\tabcolsep}{4pt}
\begin{longtable}{p{\dimexpr 0.30\textwidth-2\tabcolsep\relax}p{\dimexpr 0.60\textwidth-2\tabcolsep\relax}}
\hline
Concept & Meaning in this exercise \\
\hline
\endfirsthead
Concept & Meaning in this exercise \\
\hline
\endhead
\textbf{Automatic differentiation (AD)} & Computing derivatives by applying the chain rule to elementary operations. \\[1.5pt]
\textbf{Loss} & A scalar measure of how poorly a model fits its target. \\[1.5pt]
\textbf{Dual number} & $x+\dot x\varepsilon$, where $\varepsilon\ne0$ and $\varepsilon^2=0$. \\[1.5pt]
\textbf{Input seed} & The initial tangent or input direction chosen for a forward sweep. \\[1.5pt]
\textbf{Directional derivative} & Rate of change along $\mathbf x+t\mathbf v$; for a scalar output, $\nabla g^T\mathbf v$. \\[1.5pt]
\textbf{Jacobian matrix} & Matrix of first partial derivatives, with outputs as rows and inputs as columns. \\[1.5pt]
\textbf{Forward mode / JVP} & Propagates values and tangents to compute $J\mathbf v$. \\[1.5pt]
\textbf{Sweep / pass} & One traversal of the calculation in a given direction. \\[1.5pt]
\textbf{Computational graph} & Nodes and directed dependencies representing a calculation. \\[1.5pt]
\textbf{Adjoint} & $\bar v=\partial L/\partial v$: sensitivity of the chosen scalar output $L$ to $v$. \\[1.5pt]
\textbf{Output seed} & Initial output adjoints; $1$ for a scalar output, or weights $\mathbf w$ for several outputs. \\[1.5pt]
\textbf{Reverse mode / VJP} & Propagates adjoints backwards to compute $J^T\mathbf w$, the transpose of $\mathbf w^TJ$. \\[1.5pt]
\textbf{Accumulation} & Adding all contributions to an adjoint when a variable affects the output through several paths. \\[1.5pt]
\hline
\end{longtable}
\endgroup

\end{document}
